\section{Related work and positioning}
\label{section:related_work}

\paragraph{Sparse representation learning.}
NLDisco builds on sparse coding and dictionary learning methods that represent observations with a small number of elements from an overcomplete dictionary~\citep{olshausen_1997_sparse_overcomplete, vincent_2008_denoising_autoencoders, hoyer_2004_sparsenmf}. Closely related SAEs have been used in mechanistic interpretability to identify approximately monosemantic features in artificial neural network activations~\citep{cunningham_2023_saes, bricken_2023_towards_monosemanticity, templeton_2024_scaling_monosemanticity, dunefsky_2024_transcoders, lindsey_2024_crosscoders, ameisen_2025_circuit_tracing}. NLDisco adapts this approach to neural population recordings: each sparse latent can be related both to external experimental variables and, through its encoder weights, to the recorded units and timebins that compose it.

\paragraph{Latent models of neural population activity.}
Classical approaches such as PCA and ICA expose structure in population activity through linear components whose interpretation is guided by variance or statistical independence~\citep{hotelling_1933_pca, comon_1994_ica}. Other LVMs capture richer population geometry and temporal dynamics given a smooth state space assumption, yielding representations useful for generative and decoding tasks, although these objectives do not necessarily make individual latent dimensions interpretable~\citep{yu_2009_gpfa, pandarinath_2018_lfads, keshtkaran_2022_autolfads, schneider_2023_cebra, song_2025_langevinflow}. NLDisco instead makes the individual sparse latent the primary focus. It does not require behavioral supervision or trial structure; it reconstructs the observed population activity such that each latent's neural composition remains explicit, and complements rather than replaces trajectory-level analyses. A detailed comparison of assumptions and capabilities is provided in \autoref{table:method_comparisons}.

\paragraph{Temporal representation learning.}
Sequence models of neural population activity, including recurrent latent dynamics models and transformer-based approaches, use temporal context to improve reconstruction or prediction~\citep{sussillo_2013_rnn_dynamics, le_2022_stndt, ye_2023_ndt2, zhang_2024_mtm}. NLDisco uses temporal context for a different objective: discovering representations that emerge or evolve over time by introduction of the novel yet conceptually simple Window SED (W-SED). A W-SED treats $S > 1$ timebins as a single example and reconstructs the full window, yielding sparse latents that explicitly reconstruct timebin-by-unit activity patterns rather than just serving as directions in a latent space. Unlike standard SAEs used in language-model interpretability, which learn features at an individual sequence position, W-SEDs integrate across sequence positions.
